% Generated draft: requires mathematical and source review.
\documentclass[11pt]{article}
\usepackage[margin=25mm]{geometry}
\usepackage{fontspec}
\setmainfont{Latin Modern Roman}
\usepackage{amsmath,amssymb,mathtools,unicode-math}
\setmathfont{Latin Modern Math}
\usepackage{enumitem,needspace,xurl,hyperref}
\setlength{\parindent}{0pt}
\setlength{\parskip}{5pt}
\setlength{\emergencystretch}{3em}
\newcommand{\N}{\mathbb{N}}
\newcommand{\Z}{\mathbb{Z}}
\newcommand{\Q}{\mathbb{Q}}
\newcommand{\R}{\mathbb{R}}
\newcommand{\C}{\mathbb{C}}
\newcommand{\F}{\mathbb{F}}
\newcommand{\A}{\mathbb{A}}
\newcommand{\PP}{\mathbb{P}}
\newcommand{\E}{\mathbb{E}}
\newcommand{\eps}{\varepsilon}
\newcommand{\dd}{\,\mathrm{d}}
\begin{document}
\section*{UnsolvedMath ID: 30004417}
\textbf{Source problem code:} catalogue entry
\subsection{Definitions and mathematical statement}
The standard setting implicit in the terminology is as follows: let \(G\) be a finite group, let \(p\) be a prime, and let \(P\) be a Sylow \(p\)-subgroup of \(G\). These hypotheses are not explicitly stated in the supplied problem text. Write
\[
N_G(P)=\{g\in G:gPg^{-1}=P\},\qquad
\operatorname{Irr}_{p'}(G)=\{\chi\in\operatorname{Irr}(G):p\nmid\chi(1)\},
\]
where \(\operatorname{Irr}(G)\) denotes the irreducible complex characters of \(G\). A McKay bijection here means a bijection
\[
f:\operatorname{Irr}_{p'}(G)\longrightarrow\operatorname{Irr}_{p'}(N_G(P)).
\]
Let \(\pi:N_G(P)\to N_G(P)/P\) be the quotient map. Identify the quotient characters with their inflations:
\[
T_{G,p,P}=\{\psi\circ\pi:\psi\in\operatorname{Irr}(N_G(P)/P)\}
=\{\theta\in\operatorname{Irr}(N_G(P)):P\subseteq\ker\theta\}.
\]
Since \(N_G(P)/P\) has order prime to \(p\), this is a subset of \(\operatorname{Irr}_{p'}(N_G(P))\).

\textbf{Problem, with unresolved qualifications retained.} Which subset of \(\operatorname{Irr}_{p'}(G)\), if it exists in the intended distinguished sense, corresponds under a McKay bijection to \(T_{G,p,P}\)? What is the analogous description for groups of Lie type?

For a specified bijection \(f\), the corresponding subset is formally
\[
S_f=f^{-1}(T_{G,p,P}).
\]
This expression only identifies the inverse-image set; it does not supply the intended intrinsic characterization. The source does not specify the bijection, its required compatibility properties, or whether the desired subset should be independent of its choice. Consequently, the qualification ``if it exists'' cannot concern merely the inverse image under an already given bijection. The title suggests a canonical subset, but the required meaning of canonical is not supplied. The Lie-type clause also leaves unspecified the precise family of finite groups and the relation of \(p\) to the defining characteristic. These omissions prevent a unique, more precise formulation of the intended characterization problem.
\subsection{Short English statement}
The problem asks how to recognize the irreducible characters of degree prime to \(p\) that correspond, through the intended McKay bijection, to characters of \(N_G(P)/P\), and asks for a description for groups of Lie type. The source does not specify the intended bijection or the conditions making the subset canonical. Its open status is supplied dated curation, not an independently verified conclusion.
\subsection{Sources}
\begin{itemize}
\item Local registry: \texttt{data/unsolvedmath/problems.json}, with \texttt{catalog/problems.json} as the additional-record registry; supplied record for permanent UnsolvedMath ID \(30004417\), title ``Canonical Subsets in McKay Bijections,'' exact source problem code \texttt{catalogue entry}. Registry contents beyond the supplied record have not been independently inspected.
\item Supplied report attribution: \emph{Mini-Workshop: Kronecker, Plethysm, and Sylow Branching Coefficients and their Applications to Complexity Theory} (2021), Oberwolfach Reports, pp.\ 517--568, DOI \texttt{10.4171/owr/2020/9}. \url{https://doi.org/10.4171/owr/2020/9}. The record gives proposed year 2020 and names no individual proposer.
\item Supplied literature source: ``iMcKay--Navarro correspondence work.'' \url{https://link.springer.com/article/10.1007/s00209-026-04040-5}.
\item Supplied literature source: Giannelli, Law, Long, and Vallejo, ``Sylow branching.'' \url{https://doi.org/10.1112/blms.12584}.
\item Supplied literature source: ``Oberwolfach Report 9/2020.'' \url{https://doi.org/10.4171/owr/2020/9}.
\item Literature assessment dated exactly \texttt{2026-08-22}: the supplied curation labels the problem open and states that canonical McKay correspondences and relevant subsets are known in several solvable and Lie-type cases, while the requested general Lie-type description remains incomplete. This is supplied dated curation derived from the OWR record, not independently verified literature assessment; its time-sensitive status and cited claims require independent checking.
\end{itemize}
\end{document}
